\section{选题的背景及意义}
\label{sec:background-significance}

\subsection{选题背景}
\label{sec:research-background}

安全漏洞是系统中可被威胁源利用或触发、导致安全策略被违反的缺陷或弱点，可能损害信息的保密性、完整性或系统的可用性\cite{nist_vulnerability}。本课题主要关注软件与固件实现中的漏洞。与一般功能缺陷相比，漏洞分析不仅需要识别异常行为，还需要确定其触发条件与安全影响，程序是否崩溃并不是判断漏洞存在的唯一依据。

漏洞治理并不止于发现缺陷或发布补丁，而是涉及漏洞发现与报告、复现确认与根因分析、影响评估与优先级划分、修复或缓解措施制定、补丁验证、协调披露与发布，以及部署后的修复状态核查等环节\cite{cert_cvd_2017,sokavr,fiber}。这些环节连接漏洞发现者、软件开发者与部署维护者，既要回答漏洞在哪里、如何修复，也要确认修复是否有效、是否体现在实际运行的程序中。因此，单个环节的自动化并不能替代对后续修复结果的确认。

近年来，大语言模型的代码理解与生成能力，为漏洞治理提供了新的技术途径。通过结合代码检索、调试、编译和测试工具，智能体能够根据执行反馈调整分析过程，逐步完成漏洞调查与程序修改。Big Sleep 已展示在真实开源软件中发现未知漏洞的能力，PatchAgent 则将缺陷定位、补丁生成与验证整合到交互式修复流程中\cite{bigsleep,patchagent}。这些进展为降低漏洞分析和修复对人工操作的依赖提供了基础，也使研究重点进一步转向复杂任务条件下的分析深度与结果可靠性。

在漏洞挖掘方面，源码可用、执行环境完备且错误判据明确时的分析能力，不能直接等同于受限场景下的发现能力。对于仅能获取设备或剥离符号二进制的嵌入式场景，程序语义恢复、内部状态观察与执行环境构建均受到限制\cite{binrex,fuzzware}。协议状态违例、行为逻辑错误和资源泄漏等问题又可能不立即产生崩溃，需要结合协议、接口及资源生命周期要求判断行为是否错误，因而不能仅依靠增加测试输入或代码覆盖解决\cite{oracleproblem}。

在漏洞修复方面，异常表现位置与实际需要修改的位置可能分属不同函数，局部代码和告警信息不足以完整解释漏洞根因\cite{li2024interpvd}。上下文缺失、根因理解偏差及边界条件遗漏，也可能使模型生成局部合理但修复不充分的补丁\cite{nong2025appatch}。在修复验证方面，通过已有测试并不意味着补丁在未覆盖的输入和执行路径上仍然正确\cite{smith2015cure}；上游补丁已经发布，也不能直接证明目标二进制包含对应修复\cite{fiber}。因此，补丁生成能力的提升还需要独立的修复有效性检查与目标产物修复状态核验相配合。

基于上述背景，本课题以“基于智能体的漏洞挖掘与修复技术研究”为题，围绕“发现问题—生成修复—确认修复”的主线，开展\textbf{基于智能体的漏洞挖掘、基于智能体的漏洞自动修复、基于智能体的漏洞修复验证}三个方向的研究，分别面向受限场景中的漏洞发现、根因驱动的补丁生成，以及修复有效性与目标产物修复状态的确认，提升漏洞治理关键环节的自动化水平与可靠性。

\subsection{选题意义}
\label{sec:research-significance}

在研究层面，本课题围绕安全要求与程序行为之间的对应关系，探索大模型语义推理与程序分析、动态测试相结合的方法。通过研究行为约束如何支持漏洞判定、根因理解如何指导有效修改，以及何种程序证据能够支持修复结论，有助于深化对不完备信息条件下智能体安全分析能力及其边界的认识。同时，区分漏洞发现、修复有效性与目标修复状态的评价依据，可为相关方法的设计和比较提供更明确的标准。

在应用层面，本课题有望提高嵌入式设备与二进制组件中隐蔽安全缺陷的发现能力，降低复杂漏洞定位与补丁编写的人工成本。通过同时关注安全要求与合法行为，减少不完整修复和功能回归；通过核查目标产物中的修复实现，帮助维护者识别补丁发布与实际修复状态之间的差异。相关成果可为软件安全维护、嵌入式设备审计和第三方组件更新核查提供技术支撑。

% \section{选题的背景及意义}

% \subsection{选题背景}

% \subsubsection{漏洞攻防格局正在发生深刻变化}

% 随着软件规模持续增长与软件供应链日益复杂，安全漏洞的数量、影响面与利用速度都在快速上升。一方面，国家级漏洞库中公开漏洞数量逐年攀升，攻击者利用公开漏洞信息实施"N-day 快速武器化"的窗口期不断缩短；另一方面，大量部署中的二进制组件缺乏透明、可验证的补丁状态，使下游用户与监管方难以判断"这个漏洞到底有没有被修复"。以 Log4j（CVE-2021-44228）为代表的事件表明：单一组件缺陷可以引发整个生态的修复验证难题~\cite{log4j}。

% 与此同时，以大语言模型（LLM）为核心的\textbf{智能体（Agent）}技术正以前所未有的速度进入软件安全领域。与传统"静态规则 + 人工操作"的分析模式不同，智能体通过"模型推理 + 工具调用 + 迭代反馈"的闭环，能够自主完成代码阅读、二进制分析、测试用例生成、补丁合成等过去高度依赖专家经验的任务。Google Project Zero 与 DeepMind 的 Big Sleep 项目首次展示了 AI 智能体在真实开源项目中发现未知内存安全漏洞的能力~\cite{bigsleep}；DARPA AIxCC 大赛则系统验证了智能体在"自动发现—自动修补—自动验证"全链路中的可行性~\cite{aixcc}。\textbf{智能体正从"辅助工具"演变为漏洞攻防闭环中的核心行动者}。

% \subsubsection{漏洞攻防闭环中的三个关键环节}

% 一个完整的漏洞攻防闭环，可以抽象为三个相互衔接的核心环节：

% \begin{enumerate}
%   \item \textbf{漏洞挖掘（Discovery）}：在目标软件中发现未知漏洞，回答"哪里可能有洞"。
%   \item \textbf{漏洞修复（Automated Vulnerability Repair, AVR）}：针对已确认存在的漏洞，自动定位并合成补丁，回答"怎么把洞补上"。
%   \item \textbf{漏洞修复验证（Repair Verification）}：判定修复是否真正消除漏洞根因、目标程序是否实际包含修复，回答"洞补上了没有、补得对不对"。
% \end{enumerate}

% 这三个环节覆盖了漏洞从"未知 $\rightarrow$ 已知 $\rightarrow$ 修复"的完整生命周期，彼此共享大量基础能力：对漏洞语义的理解、对补丁意图的建模、对二进制与源码之间语义鸿沟的桥接、对修复正确性的判定。\textbf{然而，现有研究大多将三者割裂对待}，分别由不同团队、不同工具链、不同基准独立推进，导致漏洞语义、补丁语义、修复语义在各自环节中被重复建模，既浪费资源，也阻碍了跨环节能力的复用。

% \subsubsection{智能体带来的新机遇与新问题}

% 智能体技术的引入，为统一上述三个环节提供了可能：

% \begin{itemize}
%   \item \textbf{能力侧}：前沿编码智能体已能熟练调用 binutils、Ghidra、gdb 等二进制与源码分析工具，完成多步推理与决策。
%   \item \textbf{成本侧}：token 成本与推理延迟的持续下降，使"在真实仓库级代码库上大规模运行智能体"从理想变为可行。
%   \item \textbf{生态侧}：SWE-bench~\cite{swebench}、CVE-Bench~\cite{cvebench}、AIxCC~\cite{aixcc} 等评测基准的成熟，为系统化比较不同智能体安全能力提供了基础设施。
% \end{itemize}

% 但与此同时，三个关键环节在智能体视角下仍各自存在未被充分回答的核心问题：

% \begin{itemize}
%   \item 在漏洞挖掘上，传统 fuzzing 与符号执行在语义/逻辑漏洞（如认证绕过、状态机错误、配置驱动缺陷）上进展缓慢，\textbf{核心难点是"自因变量"——那些不受外部输入控制、由内部状态/配置/时序决定的变量}——导致现有输入驱动方法在错误维度上爆炸。
%   \item 在漏洞修复上，智能体已能生成"看起来合理"的补丁，但\textbf{"补丁是否真正消除了漏洞根因"}这一判定问题（即 PPT 在修复场景下的对偶问题）仍缺乏可靠机制，导致现有基准上"补丁过拟合 PoC、回归失败、绕过检查"等问题频发~\cite{rrbench}。
%   \item 在漏洞修复验证上，智能体虽已展现出与专用系统相当甚至更优的补丁存在性测试（PPT）性能~\cite{zhang2026agenticppt}，但\textbf{补丁定位（patch localization）}仍是主要瓶颈，且\textbf{对"演化补丁"（patch evolution）的鲁棒性}尚未解决——这部分能力本质上是漏洞语义理解问题。
% \end{itemize}

% 由此可见，三个环节在智能体范式下并非彼此孤立，而是被一条隐含的主线紧密连接：\textbf{对漏洞—补丁—修复三者语义的统一建模与跨环节复用}。抓住这条主线，就有可能突破现有"单点智能体应用"的局限，构建面向完整攻防闭环的智能体安全分析体系。

% \subsubsection{选题的提出}

% 基于上述观察，本文提出"\textbf{基于智能体的漏洞挖掘与修复}"这一研究课题，以"漏洞语义—补丁语义—修复语义"的统一建模为核心，系统研究智能体在漏洞挖掘、漏洞自动修复、漏洞修复验证三个关键环节中的能力边界、瓶颈机理与增强方法，并探索三个环节之间的能力迁移与反馈闭环。

% \subsection{选题意义}

% \subsubsection{理论意义}

% \begin{enumerate}
%   \item \textbf{推动"智能体 + 程序分析"交叉范式的形成}：现有智能体安全研究多停留在"把 LLM 当作更强的分类器/生成器"层面，缺乏对智能体推理过程与程序分析约束之间关系的理论刻画。本课题将探索智能体的"观察—推理—行动"循环与污点分析、符号执行、规约验证等传统程序分析技术的耦合机制，为构建新型混合分析范式提供理论支撑。
%   \item \textbf{建立漏洞语义跨环节统一表示}：针对漏洞挖掘、修复、修复验证三环节语义割裂的现状，本课题尝试建立以"漏洞触发条件—补丁干预点—修复不变式"为核心的统一表示，为后续跨任务迁移学习、多任务智能体训练提供概念框架。
%   \item \textbf{深化对"自因变量"与"演化补丁"两类共性瓶颈的理解}：前者是漏洞挖掘的核心障碍，后者是 PPT 与修复共同面对的难点。本课题将从智能体视角对二者进行系统建模，形成可被后续研究复用的分析框架。
% \end{enumerate}

% \subsubsection{实践意义}

% \begin{enumerate}
%   \item \textbf{提升漏洞挖掘的自动化水平}：相比内存破坏漏洞，语义/逻辑漏洞（认证绕过、状态机错误、配置驱动缺陷）长期缺乏高效自动化检测手段。基于智能体的挖掘方法有望通过自然语言规约理解、内部状态推断与测试用例合成，显著扩展可检测漏洞类型。
%   \item \textbf{提高自动修复的可信度}：通过将 PPT 思想引入修复验证环节，本课题旨在构建"修复—验证"闭环，减少"补丁过拟合 PoC"、"修复引入回归"等问题，提升智能体生成补丁的真实可用率，为软件供应链的快速应急响应提供工具支撑。
%   \item \textbf{支撑大规模供应链安全审计与修复状态核查}：PPT 是软件物料清单（SBOM）与漏洞影响评估的关键技术，也是修复验证能力的直接应用。基于智能体的端到端 PPT 可显著降低对符号表与人工目标函数定位的依赖，提升对剥离二进制与真实发行版包的分析能力，对关键基础设施软件资产盘点与修复状态核查具有直接价值。
%   \item \textbf{服务国家网络空间安全战略}：漏洞攻防能力是国家网络空间安全的核心竞争力之一。本课题围绕智能体这一新兴技术形态，系统研究其在漏洞攻防关键环节中的应用边界与增强路径，可为我国在下一代自动化攻防体系建设中抢占先机提供理论与技术储备。
% \end{enumerate}
